\section{معادلات الموضع الديكارتي (\LR{Cartesian Position Equations})}
\label{sec:position}

للحصول على التحويل الفراغي الكلي من قاعدة الروبوت إلى نقطة المؤثر النهائي، يتم ضرب مصفوفات التحويل المتسلسلة التي تم حسابها في القسم السابق ببعضها البعض:

\begin{equation}
T_3^0 = A_1^0 \cdot A_2^1 \cdot A_3^2
\end{equation}

وبعد إتمام عمليات الضرب المصفوفي وتطبيق متطابقات المجموع للزوايا المثلثية، تم التوصل إلى مصفوفة التحويل الكلية \LR{$T_3^0$}. ولتبسيط عرض المصفوفة، تم استخدام الاختصارات \LR{$c_{12} = \cos(\theta_1 + \theta_2)$} و \LR{$s_{12} = \sin(\theta_1 + \theta_2)$}:

\begin{equation}
T_3^0 = 
\begin{bmatrix}
c_{12} & s_{12} & 0 & L_1 c_1 + L_2 c_{12} \\
s_{12} & -c_{12} & 0 & L_1 s_1 + L_2 s_{12} \\
0 & 0 & -1 & -d_3 \\
0 & 0 & 0 & 1
\end{bmatrix}
\end{equation}

يمثل العمود الرابع من هذه المصفوفة شعاع الموضع الديكارتي للمؤثر النهائي في الفضاء التشغيلي. ومنه، يمكن استخلاص معادلات الحركة المباشرة النهائية \LR{($x, y, z$)} بدلالة الإحداثيات المفصلية \LR{($\theta_1, \theta_2, d_3$)} وأطوال الوصلات \LR{($L_1, L_2$)} كما يلي:

\begin{equation}
\begin{cases}
x = L_1 \cos(\theta_1) + L_2 \cos(\theta_1 + \theta_2) \\
y = L_1 \sin(\theta_1) + L_2 \sin(\theta_1 + \theta_2) \\
z = -d_3
\end{cases}
\label{eq:fk_position}
\end{equation}
يُلاحظ من المعادلة الأخيرة أن الإحداثي $(z)$ يأخذ قيمة سالبة $(-d_3)$، وهذا يتوافق رياضياً وفيزيائياً مع البنية الميكانيكية، حيث أن المحور $Z_2$ موجه نحو الأسفل، وبالتالي فإن الإزاحة الموجبة للمفصل الانسحابي ستؤدي إلى نزول الذراع أسفل مستوي القاعدة المرجعية.